%% =====================================================================
%%  8 Strengths and weaknesses
%%  Numbered claims, each carrying a number; aligned with FINAL_NUMBERS.md.
%% =====================================================================
\section{Strengths and Weaknesses}
\label{sec:strengths}

\subsection*{Strengths}

\begin{enumerate}[label=S\arabic*,font=\bfseries,leftmargin=*,itemsep=3pt]
  \item The answer is reproduced by an independently written pipeline: two
        parameterisations of the transverse force (a seven-parameter spin
        fit and the six-parameter perpendicular-plane fit) and two importance
        samplers (\num{40000} and \num{131072} draws) place the crossing within
        \SI{3}{mm} of each other and the goal probability within
        \num{0.002}.
  \item Model-form uncertainty is quantified adversarially rather than hidden:
        \num{25} single-factor scenarios and four joint-stress cases move the goal
        probability from \num{0.735} at \SI{20}{mm} position noise to \num{0.9997}
        at \SI{2.5}{mm}; their \SI{95}{\percent} intervals union to the scenario
        envelope $y\in[2.259,3.935]$~m, $z\in[1.169,2.698]$~m, which is reported
        separately from the narrower model envelope $y\in[2.125,3.121]$~m built
        from the calibrated intervals of the two complete models, and neither is
        called a confidence interval.
  \item The statistical intervals are calibrated, not assumed: \num{1000}
        synthetic truths give coverage \num{0.917} in $y$ and \num{0.446} in
        $z$ and the inflation factors $\lambda_y=0.96$, $\lambda_z=4.94$,
        while the profile and Wald widths agree to \num{0.1} and \num{0.6} per
        cent; a \num{500}-experiment validation covers the truth \num{0.918}
        in $y$, \num{0.928} in $z$ and \num{0.852} jointly.
  \item The model order is justified by the data rather than chosen: the
        seven-parameter Jacobian is $6\times7$ of rank \num{6} with singular
        values from \num{327} to \num{1.75}, prior shrinkages of \num{8.69},
        \num{69.99} and \num{72.65} per cent imply \num{5.27} data-determined
        parameters, so the six-parameter fit is the largest model the fixes
        can support without a prior.
  \item The method is real-time and comes with a design: fitting costs
        \SI{8}{ms} plus \SI{0.5}{ms}, so the prediction is ready
        \SI{0.11}{s} after the kick, \SI{135}{ms} before a keeper can resolve
        height; at the adopted \SI{5}{mm} noise the problem's timings give
        $\sigma_y=\SI{211}{mm}$ while the best pair, \SI{0.11}{s} and
        \SI{0.31}{s}, gives \SI{12}{mm}, a factor \num{18}.
\end{enumerate}

\subsection*{Weaknesses}

\begin{enumerate}[label=W\arabic*,font=\bfseries,leftmargin=*,itemsep=3pt]
  \item Two fixes leave no redundancy, so a model with six free parameters has
        no goodness-of-fit test: that is exactly the degenerate case M0 of
        Table~\ref{tab:ladder}, whose residual is identically zero and whose
        agreement with the data is therefore not evidence. The adopted model has
        six parameters for six numbers only in the same counting sense; its
        every fit is penalised by the priors, so the MAP residual
        (\SI{1.45}{mm}) is not zero, but the model still has no spare degree of
        freedom. Only the nested physical models can misfit, with rms
        \SI{5.03}{mm} without a lateral force against \SI{0.54}{mm} with one.
  \item The position noise is assumed, not measured: $\sigma_{\text{pos}} =
        \SI{5}{mm}$ is half of the last printed digit of the given
        coordinates, and it drives both the point and the probability---at
        \SI{20}{mm} the goal probability falls to \num{0.735} and the scenario
        means move $y$ over \SIrange{2.687}{3.130}{m}, and at the
        \SIrange{55}{65}{mm} noise of open tracking data the fit is
        prior-dominated.
  \item The transverse coefficients are effective, not identified: the fit
        returns $C_D=\num{0.347}$ (\num{0.378} in the prior-free
        decomposition) against \num{0.19} to \num{0.22} measured in wind
        tunnels at these speeds \citep{carre2005,goff2012}, so part of the transverse effect may be an
        unmodelled along-track or wake force; the spin is never observed, and
        the implied spin rate is a consequence of the assumed lift law rather
        than a measurement.
  \item The end-to-end validation is synthetic: no public three-dimensional
        trajectory set with measured spin and a known crossing point exists, so
        the coverage results test the estimator and not the physics, while real
        tracking data enter only as measurement precision (median \SI{0.050}{m}
        for broadcast 3-D reconstruction, \SI{5.4}{}--\SI{6.5}{cm} for open 2-D
        tracking); the scenarios
        vary one assumption at a time apart from four joint-stress cases, which
        leaves correlated model error only partly explored, and the adopted-model
        coverage rests on fifty and twenty-five trials, so it detects only gross
        failures; wind is not separable from the launch velocity with two fixes,
        its \SI{0.007}{m} refit effect being a property of the fit rather than of
        the atmosphere.
\end{enumerate}
